
\section{Methodology}
\label{sec:methodology}

\subsection{Research Design}
\label{subsec:research-design}

This study adopts a controlled experimental approach to
investigate how evidence position and context length affect
information retrieval accuracy in large language models (LLMs).

The primary experiment uses a factorial design with two
independent variables: the relative position of relevant
evidence and the total context length. A synthetic key--value
retrieval task is used to control the experimental conditions
and minimize variation in retrieval difficulty.

The experimental framework is expressed as
\begin{equation}
    Q = f(p,L \mid m,d_0),
    \label{eq:primary-framework}
\end{equation}
where $Q$ represents retrieval accuracy, $p$ is the relative
position of relevant evidence, $L$ is the total context
length, $m$ identifies the evaluated model, and $d_0$
represents a fixed baseline distractor generation policy.

Under $d_0$, distractor records are generated using a
consistent sampling procedure without deliberately
introducing semantic similarity, conflicting evidence,
or repeated target keys. The policy remains unchanged
across experimental conditions, although the number
of distractor records increases with context length.

Model parameters remain fixed throughout the experiments.
Results are analyzed separately for each model.

\subsection{Experimental Variables}
\label{subsec:variables}

\subsubsection{Independent Variables}

Two independent variables are manipulated.

\begin{itemize}
    \item \textbf{Evidence Position ($p$):}
    The target record is inserted at five relative positions:
    0\%, 25\%, 50\%, 75\%, and 100\%.
    These positions are defined relative to the sequence
    of key--value records, excluding the fixed instructions
    and retrieval question.

    \item \textbf{Context Length ($L$):}
    Four target context lengths are considered:
    2K, 4K, 8K, and 16K tokens.
    The actual token counts are verified using the tokenizer
    associated with each evaluated model.
\end{itemize}

\subsubsection{Dependent Variable}

The dependent variable is retrieval correctness, measured
as a binary outcome for each test case:

\begin{equation}
    Y_i =
    \begin{cases}
        1, & \text{if the answer is correct},\\
        0, & \text{otherwise}.
    \end{cases}
    \label{eq:correctness}
\end{equation}

Retrieval accuracy is calculated by aggregating correctness
across test cases within each experimental condition.

\subsubsection{Controlled Variables}

To isolate the effects of position and context length,
the following conditions are controlled:

\begin{itemize}
    \item The target key--value pair and expected answer
    within each matched test case.
    \item The retrieval question and instruction format.
    \item The key and value character sets and lengths.
    \item The structure and formatting of all records.
    \item The baseline distractor generation policy $d_0$.
    \item Model parameters and inference settings.
    \item The output format and evaluation criteria.
\end{itemize}

\subsection{Dataset and Context Construction}
\label{subsec:dataset}

A synthetic key--value retrieval dataset is generated
to provide precise control over evidence position
and context length.

Each record contains a key and an associated value.
Keys are unique within each context, and the target
key appears exactly once. Keys and values are generated
using predefined character sets and fixed character
lengths to standardize record structure.

For example, a simplified context may contain:

\begin{verbatim}
Key: A7F92
Value: BLUE

Key: B3K81
Value: GREEN

Key: C9M24
Value: RED

Key: D2P65
Value: WHITE

Key: E8R13
Value: BLACK
\end{verbatim}

The corresponding retrieval question is:

\begin{quote}
\textit{What is the value associated with key C9M24?
Answer with the value only.}
\end{quote}

The expected answer is \texttt{RED}.

In this example, the record containing
\texttt{C9M24} is the target evidence,
while the remaining records serve as distractors.

For each context length, the same set of distractor
records is reused across positional conditions.
The target record is inserted at different relative
positions while preserving the relative order of
distractor records.

As context length increases, additional distractor
records are generated using the same baseline policy
$d_0$. Thus, $d_0$ specifies a consistent generation
procedure rather than a fixed number or set of records.

Although key and value character lengths are fixed,
their tokenized lengths may differ. Actual token counts
are therefore measured, and the resulting context
lengths and evidence positions are recorded.

\subsection{Experimental Procedure}
\label{subsec:procedure}

The experiment follows five steps.

\begin{enumerate}
    \item \textbf{Generate retrieval instances.}

    Construct a collection of target key--value pairs
    with unique identifiers and known correct answers.

    \item \textbf{Construct controlled contexts.}

    For each retrieval instance and context length,
    generate a baseline set of distractor records.
    Insert the same target record at each selected
    relative position.

    \item \textbf{Execute model inference.}

    Submit each context and its retrieval question
    to the selected model. Use deterministic decoding
    where supported, with a fixed output-token limit.

    \item \textbf{Evaluate responses.}

    Compare the model response against the expected
    value using exact-match evaluation. Invalid,
    missing, or incorrect answers are counted
    as retrieval failures.

    \item \textbf{Repeat across conditions.}

    Repeat the procedure over multiple independently
    generated retrieval instances. Record the model,
    context length, evidence position, expected answer,
    predicted answer, and correctness for each trial.
\end{enumerate}

\subsection{Evaluation Metrics}
\label{subsec:evaluation}

The primary evaluation metric is retrieval accuracy:

\begin{equation}
    \operatorname{Accuracy}(p,L)
    =
    \frac{N_{\mathrm{correct}}(p,L)}
         {N_{\mathrm{total}}(p,L)}.
    \label{eq:accuracy}
\end{equation}

To characterize positional differences, the study
also measures the positional accuracy gap:

\begin{equation}
    \Delta_{\mathrm{pos}}(L)
    =
    \max_p \operatorname{Accuracy}(p,L)
    -
    \min_p \operatorname{Accuracy}(p,L).
    \label{eq:position-gap}
\end{equation}

This metric quantifies performance variation across
evidence positions without assuming that all models
exhibit a U-shaped accuracy pattern.

The interaction between position and context length
is examined by comparing positional accuracy curves
and positional gaps across different context lengths.

Results are visualized using accuracy-versus-position
plots, with separate curves for each context length.
Where the number of trials permits, confidence
intervals are reported to indicate measurement
uncertainty.

\subsection{Reproducibility and Limitations}
\label{subsec:limitations}

The experiment records dataset-generation seeds,
model versions, inference configurations, token counts,
prompts, and evaluation results to support reproducibility.

Synthetic key--value retrieval provides a controlled
environment for investigating evidence position and
context length. However, it primarily evaluates
exact information retrieval rather than complex
semantic reasoning.

Furthermore, increasing context length requires
additional distractor records under the fixed record
structure. Consequently, context length and the
number of distractor records are not independently
manipulated in this design.

The observed length effects therefore reflect
increasing context size under a fixed distractor
generation policy, rather than the isolated effect
of token length alone.

The findings will be interpreted as evidence about
controlled information retrieval and will not be
generalized directly to all forms of long-context
reasoning or agentic tasks.
